\section{Interpreting population SD and conditional SD in the examples}
\label{sec:sd-interpretation}

\subsection{Which experiment does each SD describe?}
The \emph{population SD} of the actual mean describes repeated draws of a
whole \(N=r+k\)-test battery from the stated population model. Every
score, including the first \(r\), is random in this experiment. On the
IQ scale the mean of those battery averages is \(100\), and their SD is
\(15\sqrt{d_N/N}\). The population SD column is consequently a property
of a battery size and its covariance structure, rather than a spread around
the particular fixed starting IQ of \(150\).

The \emph{conditional SD} describes the remaining variation in actual
future outcomes after the first \(r\) scores have been observed to equal
\(150\). Those first scores contribute fixed values to the inclusive
mean; the later scores remain random and correlated under the conditional
Gaussian law. The centre of this conditional distribution is
\(m_{r,k}^{\mathrm{IQ}}\), and its SD is
\(15\sqrt{k(1-c)d_N/d_r}/N\), by
Theorem~\ref{thm:multiple-sd} and Corollary~\ref{cor:iq-reporting}.
The conditional SD is independent of the numerical starting value in this
Gaussian model, although the conditional mean depends on that value.

Replacing every later outcome by its conditional prediction produces a
third object: a fixed list of numbers. Its recursive average is the
conditional centre just described and has SD \(0\). The nonzero SD
columns describe actual outcome means, not a repeated computation of that
fixed list. These distinctions follow from the moment and conditioning
theorems in Section~\ref{sec:foundations}.

For example, two initial IQ-150 scores and ten subsequent predictions give
an inclusive conditional mean of \(1675/12\approx139.583333\).
For a randomly drawn twelve-test battery, the population SD is
\(\sqrt{285/2}\approx11.937336\) IQ points. After the first two scores
are fixed at \(150\), the actual twelve-score mean has conditional SD
\(\sqrt{475/16}\approx5.448624\) IQ points. The recursively filled list
\(150,150,137.5,\ldots,137.5\) has the fixed arithmetic mean
\(1675/12\) and SD zero for that mean. Section~\ref{sec:certificates}
proves the displayed decimal rounding.

\subsection{The variance decomposition}
\begin{proposition}[The contribution of the observed starting block]
\label{prop:variance-components-general}
For \(0\leq c<1\), let \(Z=(X_0,\ldots,X_{r-1})^{\T}\) be a random
starting block, without constraining its entries to be equal. Then
\begin{align}
 \E[C_{r,k}\given Z]
  &=\frac{d_N}{N d_r}\sum_{i=0}^{r-1}X_i,
    \label{eq:random-starting-block-centre}\\
 \Var\bigl(\E[C_{r,k}\given Z]\bigr)
  &=\frac{r d_N^2}{N^2d_r},
    \label{eq:starting-block-centre-variance}\\
 \Var(C_{r,k})
  &=\frac{r d_N^2}{N^2d_r}
    +\frac{k(1-c)d_N}{N^2d_r}
   =\frac{d_N}{N}.
    \label{eq:general-total-variance}
\end{align}
\end{proposition}
\begin{proof}
Theorem~\ref{thm:regression} gives each future conditional mean as
\((c/d_r)\sum_{i<r}X_i\). Conditional linearity, with the first block
known, therefore gives
\[
 \E[C_{r,k}\given Z]
 =\frac1N\left(1+\frac{kc}{d_r}\right)\sum_{i<r}X_i
 =\frac{d_N}{N d_r}\sum_{i<r}X_i.
\]
Theorem~\ref{thm:moments} gives
\(\Var(\sum_{i<r}X_i)=r d_r\); multiplying by the square of the
displayed coefficient proves
\eqref{eq:starting-block-centre-variance}. The conditional variance from
Theorem~\ref{thm:multiple-sd} is independent of the particular block
values. The law of total variance in Theorem~\ref{thm:conditional}
therefore gives the sum in \eqref{eq:general-total-variance}.
Finally,
\(r d_N+k(1-c)=N d_r\), by expanding with \(N=r+k\), which simplifies
the sum to \(d_N/N\).
\end{proof}

Thus population variance contains both variation of conditional centres
across starting blocks and the variation remaining within a conditioned
model. In the two-start, ten-prediction IQ example, this exact decomposition
is
\[
 \frac{285}{2}=\frac{1805}{16}+\frac{475}{16}
 \quad\text{IQ points squared}.
\]
The decomposition adds variances. The square roots of the total variance
and of the remaining conditional variance are the two SDs reported in the
example table. It is a population statement about random starting
blocks; a separate random distribution for several initial scores forced
to share a single common value has not been postulated.

\subsection{Why the conditional SD can increase with more predictions}
At \(c=0.6\), the conditional SD columns for \(k=1,2,10\) increase
towards the positive limits in Corollary~\ref{cor:iq150-limits}, even
while the population SD decreases. Initially, the fixed starting scores
account for a large fraction of the mean, and their conditional variance
is zero. As more future tests enter the average, that fixed fraction
decreases. Correlation leaves a common uncertain component in the future
mean, which accounts for a positive limiting conditional SD when
\(0<c<1\).

\begin{proposition}[The directions of change in the selected examples]
\label{prop:conditional-sd-growth-examples}
For every fixed \(r\geq1\) at \(c=3/5\), the population variance of
\(C_{r,k}\) strictly decreases with integer \(k\geq1\), while its
conditional variance given the initial block strictly increases.
\end{proposition}
\begin{proof}
The population variance is \(c+(1-c)/(r+k)\), so increasing \(k\) by
one changes it by
\(- (1-c)/[(r+k)(r+k+1)]<0\).
For the conditional variance put
\[
 f(k)=\frac{k(d_r+ck)}{(r+k)^2},\qquad
 \Var(C_{r,k}\given Z)=\frac{1-c}{d_r}f(k).
\]
Direct multiplication gives
\[
 f(k+1)-f(k)
 =\frac{(2cr-d_r)k^2+(2cr^2+2cr-d_r)k+r^2(c+d_r)}
        {(r+k)^2(r+k+1)^2}.
\]
At \(c=3/5\), the three numerator coefficients are
\[
 \frac{3r-2}{5},\qquad
 \frac{6r^2+3r-2}{5},\qquad
 \frac{r^2(3r+5)}{5},
\]
which are all strictly positive for \(r\geq1\). The denominator and
\((1-c)/d_r\) are positive as well. This proves the increasing direction.
Taking square roots preserves both strict comparisons, so the same
directions apply to SDs. Multiplication by \(15\) gives the IQ-scale
statements.
\end{proof}
This direction of conditional change is established for the selected
correlation; no universal monotonicity claim for every admissible \(c\)
is needed.

\subsection{What population and conditional simulations should sample}
For the nonnegative-correlation examples, one exact aggregate Gaussian
representation of the population mean is
\begin{equation}
 C_{r,k}\ \stackrel{d}{=}
 \sqrt c\,G+\sqrt{\frac{1-c}{N}}\,E,
 \label{eq:population-aggregate-sampler}
\end{equation}
where \(G,E\) are independent standard normals. A corresponding
conditional representation is
\begin{equation}
 C_{r,k}\given Z=a\one_r\ \stackrel{d}{=}
 B_{r,k}+\frac1N\left[
   k\sqrt{\frac{c(1-c)}{d_r}}\,G+
   \sqrt{k(1-c)}\,E\right].
 \label{eq:conditional-aggregate-sampler}
\end{equation}
Here \(\stackrel{d}{=}\) denotes equality in distribution, rather than
equality of individual draws. The Gaussian affine-closure and uniqueness
theorem, Theorem~\ref{thm:gaussian}, proves these representations: their
means and variances, computed using Theorem~\ref{thm:moments}, are exactly
those in Theorem~\ref{thm:multiple-sd}. On the IQ scale, apply
\(100+15C_{r,k}\) to either representation.
At \(c=1\), the conditional representation reduces to the constant
\(a\), in agreement with the equal-score support argument in
Proposition~\ref{prop:multiple-boundaries}.

A population simulation lets the starting tests vary together with the
other tests, while a conditional simulation samples the conditional law
after the starting scores are fixed. The latter can be generated directly
from \eqref{eq:conditional-aggregate-sampler}; exact equality to 150 is a
probability-zero event in the unconditioned continuous model, so accepting
unconditioned draws only when the starting scores equal 150 is not a
conditional sampler. Empirical means and SDs from finite simulation runs
are numerical diagnostics of the implemented model and generator. The
mathematical calculations and proofs above supply the target values.
